\documentclass[11pt]{article}

\usepackage[margin=1in]{geometry}
\usepackage[T1]{fontenc}
\usepackage[utf8]{inputenc}
\usepackage{lmodern}
\usepackage{xcolor}
\usepackage{xurl}
\usepackage{parskip}

\newcommand{\missing}[1]{\textcolor{red}{\textbf{[MISSING: #1]}}}

\pagestyle{empty}

\begin{document}

\noindent
Faculdade de Sa\'ude P\'ublica, Universidade de S\~ao Paulo, S\~ao Paulo, Brazil \\
\missing{submission date}

\vspace{1em}
\noindent
The Editors, \emph{International Journal of Forecasting}

\vspace{1em}

\noindent Dear Editors,

\vspace{0.6em}

We submit for your consideration ``Series length, not model choice, drives forecasting
accuracy of cardiovascular mortality in S\~ao Paulo, Brazil: a benchmark of foundation,
classical and boosting models with measured uncertainty'', as an original research article.

Once uncertainty is measured, a zero-shot foundation model (TimesFM) and two classical
models (SARIMA, Prophet) are statistically indistinguishable for forecasting monthly
cardiovascular mortality, and what moved accuracy was the length of the training series
rather than the choice of model. The spread among the three leading models is 0.13
percentage points against a mean 95\% interval width of 1.07, while nine additional years
of training reduced error by 1.23 to 1.71 points on identical test dates.

The work fits your scope of forecasting mortality and forecasting for healthcare, and it is
the kind of evaluation study the journal emphasises: 618 out-of-sample forecasts per model
from 103 rolling origin windows, with every comparison carrying a block bootstrap interval
over whole windows and a Diebold-Mariano test at each horizon, under a decision criterion
fixed before the analysis.

It is also, in part, a replication of our own earlier round on a shorter series, in which
TimesFM ranked first by 0.38 percentage points of sMAPE. That ranking survived neither nine
further years of data nor the addition of intervals, and the ordering reversed. We report
this explicitly because the lesson generalises past our series: gaps of a few tenths of a
point are routinely published as rankings in applied health forecasting, while the interval
around aggregate error over a few hundred dependent forecasts is on the order of a full
point.

On your data and code policy and the reproducibility check, the aggregated series, every
forecast file and every uncertainty output are already public at
\url{https://github.com/fabianofilho/cardiovascular-timeseries-prediction}. Every table and
figure is generated programmatically from the forecast files by one script, which also
emits a machine-readable file of every value cited in the text; none was transcribed by
hand. We verified that a clean clone reproduces the manuscript, including the compiled PDF,
with no local data. Individual death records are not redistributed, as the included
extraction script retrieves them from the public Ministry of Health source.

The manuscript is original, unpublished, and not under consideration elsewhere. All authors
approved the submission and declare no competing interests. As declared in the manuscript,
an AI coding assistant was used for implementation, analysis scripting and drafting; it is
not an author, and the authors verified all analyses and values and take responsibility for
the content.

\missing{suggested reviewers, if the submission system requests them; funding statement to
match the manuscript declarations}

\vspace{0.8em}

\noindent Sincerely, \\[0.6em]
Fabiano B. N. Filho, on behalf of the authors \\
\missing{corresponding author email and ORCID}

\end{document}
